%% 20_dscircuits.tex -- Chapter 20: $\Delta\Sigma$ Circuits: Discrete-Time, Continuous-Time and Incremental
%% Status: OUTLINE (build stage 4). Replace the outline with the full chapter when it is written.
\chapter{$\Delta\Sigma$ Circuits: Discrete-Time, Continuous-Time and Incremental}\label{ch:dscircuits}
\chapterlead{From the loop filter on paper to integrators in silicon, and to the incremental converter that serves the high-resolution \OB{} channels.}

\begin{specnote}[Chapter status: outline]
This chapter is an outline. It will be written in build stage 4 (the rest of Part III with Labs 7, 8 and 10). The topics, the design question and the labs below are fixed; the derivations, worked examples and figures are still to come.
\end{specnote}

\begin{objectives}
\begin{itemize}
  \item build switched-capacitor integrators and set their OTA requirements;
  \item design continuous-time loops with resistive or Gm-C integrators and use their inherent anti-aliasing;
  \item handle excess loop delay;
  \item design an incremental $\Delta\Sigma$ converter for multiplexed channels;
  \item design and budget the decimation filter (CIC/sinc$^K$ and half-band FIR).
\end{itemize}
\end{objectives}

\begin{insight}[Design question]
How do you implement the 2nd-order incremental modulator of \OB{} and its decimator within the power budget?
\end{insight}

\section{Switched-capacitor integrators}\label{sec:ch20-1}
Charge transfer, OTA gain and slew requirements, and $kT/C$ noise of the sampling capacitor.

\section{Continuous-time \texorpdfstring{$\Delta\Sigma$}{Delta-Sigma}}\label{sec:ch20-2}
Resistive and Gm-C integrators; the implicit anti-aliasing of a CT loop; clock-jitter sensitivity.

\section{Excess loop delay}\label{sec:ch20-3}
How a fraction of a clock period of delay destabilises the loop and how to compensate it.

\section{Incremental \texorpdfstring{$\Delta\Sigma$}{Delta-Sigma}}\label{sec:ch20-4}
Resetting the loop for every sample so channels can be multiplexed; the resolution versus number of cycles.

\section{Decimation}\label{sec:ch20-5}
CIC (sinc$^K$) and half-band FIR stages, their word lengths and their power.

\section{The reference design}\label{sec:ch20-6}
The full 2nd-order incremental modulator of the \OB{} ECG/EEG mode.

\begin{tryit}[Labs for this chapter]
Lab~10 (\cref{sec:lab10}) implements the Lab~9 loop with switched-capacitor integrators and a 1-bit quantiser in \sgx{} and compares it with the ideal model.
\end{tryit}

\paragraph{Read with this chapter} Schreier, Temes and Pavan~\cite{schreier}.
